% Proposal: a new section between Results and Discussion, for the person's decision (docs/seats/WRITER.md:
% until a venue is chosen, new text reaches the master only as a proposal). Asked by the person, 2026-11-18:
% "add a section on agentic loop usage from longevity-loci repo ... describe agentic loop structure and
% interaction mechanisms as a showcase of autonomous AI loop + module creation + HITL". The orchestration kit
% the loop runs on is not named, on the person's word.
%
% Facts read from the longevity-loci repository at origin/main 662e366 (2026-11-18): docs/seats/*.md,
% docs/DECISIONS.md, rows/PROGRESS.md, docs/seats/MODULE_NOTES.md (the 2026-11-16 rehearsal),
% docs/longevitymap_v2/longevitymap_v2_requirements.md, git log trailers.
%
% Needs in references.bib (checked at PubMed 2026-11-18, PMID 23998809):
% @article{budovsky2013longevitymap,
%   title   = {{LongevityMap}: a database of human genetic variants associated with longevity},
%   author  = {Budovsky, Arie and Craig, Thomas and Wang, Jingwei and Tacutu, Robi and Csordas, Attila and
%              Louren{\c{c}}o, Joana and Fraifeld, Vadim E. and de Magalh{\~a}es, Jo{\~a}o Pedro},
%   journal = {Trends in Genetics}, volume = {29}, number = {10}, pages = {559--560}, year = {2013},
%   doi     = {10.1016/j.tig.2013.08.003}}

% ================================================================
\section{Case study: an autonomous curation loop with human checkpoints}
\label{sec:loop}
% ================================================================

The sections above follow one assistant authoring one module with its author
at the keyboard. The same tools also serve a longer-running arrangement, in
which several agent sessions work on one module for days while the maintainer
is mostly away. We ran such a loop from 4 to 18 November 2026 to rebuild a
longevity module from its evidence. The module's first version is a port of
HAGR's LongevityMap~\citep{budovsky2013longevitymap}, which has not changed
since 2017. Most of its evidence is from candidate-gene studies, and more than
half of its study-cited variants rest on a single centenarian classifier whose
predecessor paper was retracted. The task was to re-read every locus against
its studies, add the genome-wide literature published since, and keep, re-tier
or drop each row according to what the papers show.

\begin{figure}[ht]
\centering
\begin{tikzpicture}[
  role/.style={draw, rounded corners, fill=gray!8, minimum width=1.7cm,
               minimum height=0.75cm, font=\small, align=center},
  art/.style={draw, dashed, minimum width=1.3cm, minimum height=0.6cm,
              font=\footnotesize, align=center},
  hum/.style={draw, thick, rounded corners, fill=yellow!15, minimum width=1.7cm,
              minimum height=0.75cm, font=\small, align=center},
  arr/.style={-{Stealth[length=2mm]}, thin},
  lab/.style={font=\scriptsize, inner sep=1pt},
  node distance=0.9cm and 1.25cm]
\node[role] (scout) {scout};
\node[art, right=of scout] (queue) {queue};
\node[role, right=of queue] (curator) {curator};
\node[art, right=of curator] (rows) {rows +\\notes};
\node[role, right=of rows] (module) {module\\builder};
\node[role, above=0.6cm of module] (liaison) {liaison};
\node[role, below=of curator] (lead) {coordinating\\agent};
\node[art, left=of lead] (dec) {decision\\record};
\node[art, below=of module] (test) {test\\registry};
\node[art, below=of lead] (desk) {question\\queue};
\node[hum, right=1.5cm of desk] (human) {maintainer};
\node[art, right=1.5cm of human] (prod) {production\\catalog};
\draw[arr] (scout) -- (queue);
\draw[arr] (queue) -- node[lab, above]{wakes} (curator);
\draw[arr] (curator) -- (rows);
\draw[arr] (rows) -- node[lab, above]{wakes} (module);
\draw[arr] (rows) |- node[lab, pos=0.3, right]{batch review} (lead);
\draw[arr] (lead) -- node[lab, above]{records} (dec);
\draw[arr] (lead) -- node[lab, right]{asks} (desk);
\draw[arr] (human) -- node[lab, above]{answers} (desk);
\draw[arr] (module) -- node[lab, right]{rehearses} (test);
\draw[arr] (module) -- node[lab, right]{tool refusals} (liaison);
\draw[arr, thick] (human) -- node[lab, above]{publishes} (prod);
\end{tikzpicture}
\caption{The curation loop. Agent roles (grey) hand work to each other through
files under version control (dashed). A push to a watched file wakes the role
that reads it. The maintainer (yellow) answers queued questions and alone takes
the irreversible steps; every other decision is written into the decision
record, together with the alternative it rejected.}
\label{fig:loop}
\end{figure}

\paragraph{Structure.} The loop consists of standing roles. Each role is held
by one long-running agent session that has a written brief, a list of the
files it may change, and a statement of what it may push
(Figure~\ref{fig:loop}). The \emph{scout} searches the recent literature track
by track. It records every query with its date window and hit counts so that
the next sweep reruns the same query, merges a paper found through several
sources into one entry, and queues each paper that bears on an existing locus,
together with the passage that shows it. The \emph{curator} works locus by
locus and reads each article at the deepest level that decides the evidence
tier (full text over abstract, the supplement for per-variant numbers). For
every article it writes a note, which is never rewritten, only extended, and
which carries the verbatim quote each conclusion rests on and the depth at
which it was read. It then writes the locus row with its tier and evidence and
commits the row and its notes together. The \emph{module builder} exports the
curated rows into a module and runs this plugin's validate, enrich and strict
compile steps on each increment. It then rehearses the version on the test
registry. A \emph{coordinating agent} writes the repository's own checks
(every citation resolves to a note, and every quote is found in the article's
cached text or is reported as unknown), reviews the curator's rows in batches
against the papers, and settles questions of how a rule should be read. Two
support roles complete the set: a \emph{runner} submits bulk article fetches
as batch jobs on a compute cluster, and a \emph{liaison} tracks releases of
the plugin, compiler and registry and files defects with whichever project
owns them.

\paragraph{Interaction mechanisms.} No person relays work between the roles.
Three mechanisms take that place. First, \emph{state lives in documents}:
rows, notes, the queue, the progress list and each role's own notes are files
in one repository. A message only delivers news, and a message that reports
work names the commit that holds it. Before a session resets, the role
rewrites its notes to describe the current state, so the next session
continues from the files rather than from a transcript. Second, \emph{wake-ups
are events}. A watcher on a path wakes the role that consumes it whenever
another role pushes to that path. The scout's queue wakes the curator, and new
rows wake the module builder. A direct message from a named role also wakes
its addressee, while other traffic waits for the next turn. Third, \emph{write
surfaces are narrow}. The curator never edits the module specification, and
the module builder never edits a row. When a row does not compile, the module
builder sends it back to the curator with the locus and field named, and a
failure on the tool's side goes to the liaison with its reproduction. The
plugin's own rules (Section~3.3) bind every role that calls it: a PMID comes
from a search result, a hand edit is logged, and a check that could not run is
reported as unknown.

\paragraph{Where the human sits.} The maintainer decides; the maintainer does
not operate the loop. Four checkpoints connect the maintainer to it.
(i)~A \emph{question queue}. A question that the written requirements leave
to the module's owner becomes an issue that can be answered from a phone. Each
issue lists its options with their costs and a default, and work that does
not depend on the answer continues meanwhile. Two scope questions were
answered in this way: APOE is kept in this module and scored once, with one
result per $\varepsilon$ pair, and biological-ageing markers form an unscored
section. (ii)~A \emph{standing delegation}. Readings of the evidence rules
that the requirements left open were settled by the coordinating agent under
the maintainer's standing permission to proceed autonomously. Examples are
what counts as a replication, when a non-replication becomes a red flag, and
how a multivariate test supports a row. Each reading is recorded as a numbered
decision together with the alternative it rejected, so that the maintainer can
review and reverse it later. The repository holds 35 such decisions for these
two weeks. (iii)~\emph{Irreversible steps are reserved}. Only the maintainer
may publish to the production catalog or close the module. Every version
produced so far is a rehearsal on the test registry. (iv)~A \emph{shared
budget}. All roles run on one account and pause at 95\% of its usage window,
so that capacity remains for the maintainer's own work.

\paragraph{What it produced.} The rehearsal of 16 November held 164 rows. These
cover all 161 seed loci, which collapse into 156 rows, plus eight secondary
signals. Their tiers are A~15, B~73, C~41 and D~35, where D is the
negative-evidence layer that records a former claim that did not hold. The
rows comprise 489 variant rows and 679 study rows. All 245 rsIDs resolved,
with no reference-allele mismatch, and the module compiled under strict
resolution. By 18 November the corpus held 140 article notes, and the
repository held 1,171 commits, 674 of them the curator's.

\paragraph{Limits.} This is one loop on one module, run with one model family.
The rows have not been compared with an independent human curation, and
production waits on the maintainer's review. The tiers therefore record what
the agents read and how the written rules classify it. They are not an
expert's verdict. The recorded decisions make that review feasible at this
scale: a reviewer can audit the rules the loop applied, rather than re-derive
every individual row.
